% =====================================================================
%  PHẦN MỞ ĐẦU
% =====================================================================
\chapter*{Phần Mở Đầu}
\addcontentsline{toc}{chapter}{Phần Mở Đầu}
\markboth{Phần Mở Đầu}{}

\section*{1. Bối cảnh và tính cấp thiết của đề tài}

Trí tuệ nhân tạo (Artificial Intelligence -- AI) và học máy (Machine Learning -- ML) đã rời phòng thí nghiệm để
trở thành một thành phần thường trực của phần mềm nghiệp vụ: hệ thống sàng lọc bệnh, công cụ định giá tài sản,
bộ lọc phản hồi khách hàng, ứng dụng nhận dạng ảnh, mô hình dự báo tài chính. Điểm chung của các hệ thống này là
chúng không được lập trình bằng luật viết tay mà \emph{học từ dữ liệu}. Sự chuyển dịch đó gắn với ba điều kiện
xuất hiện cùng lúc trong thập niên 2010: dữ liệu lớn, phần cứng tính toán song song (GPU) và các thuật toán tối ưu
cho mạng nhiều tầng \cite{goodfellow2016}.

Một kỹ sư phát triển hệ thống thông minh hiếm khi chỉ gặp một dạng dữ liệu. Trong thực tế có ít nhất ba hình thái,
mỗi hình thái hợp với một họ mô hình riêng:

\begin{itemize}
  \item \textbf{Dữ liệu bảng} (y tế, bất động sản, tài chính doanh nghiệp): mỗi dòng là một đối tượng độc lập với
    các thuộc tính số và phân loại. Đây là địa hạt của các thuật toán học máy cổ điển (hồi quy tuyến tính,
    Logistic Regression, cây quyết định, Random Forest, SVM) và của mạng nơ-ron đa tầng (MLP).
  \item \textbf{Dữ liệu không gian} (ảnh): các điểm ảnh kề nhau tương quan chặt, và cùng một mẫu hình có thể xuất
    hiện ở bất kỳ vị trí nào. Mạng nơ-ron tích chập (CNN) khai thác đúng hai tính chất này nhờ chia sẻ trọng số
    và vùng tiếp nhận cục bộ.
  \item \textbf{Dữ liệu chuỗi} (giá chứng khoán, nhật ký hành vi người dùng): thứ tự thời gian mang thông tin, các
    mẫu không còn độc lập. Mạng nơ-ron hồi quy (RNN) và các biến thể có cổng (LSTM, GRU, BiLSTM) mang một trạng
    thái ẩn đi qua từng bước thời gian để ghi nhớ ngữ cảnh.
\end{itemize}

Từ yêu cầu đó, tiểu luận \emph{“Nghiên cứu và triển khai các mô hình học máy cơ bản, mạng nơ-ron tích chập (CNN)
và mạng nơ-ron hồi quy (RNN) trong hệ thống thông minh”} đặt mục tiêu đi trọn một vòng từ lý thuyết đến sản phẩm
cho cả ba hình thái dữ liệu: hiểu toán học của mô hình, tự cài đặt lại nó, huấn luyện trên dữ liệu thật, đo đạc
khách quan rồi đóng gói thành ứng dụng web chạy được.

\section*{2. Mục tiêu và nhiệm vụ}

Mục tiêu tổng quát là xây dựng nền tảng kiến thức vững cùng năng lực thực hành trên toàn bộ tiến trình từ học máy
cổ điển đến các kiến trúc học sâu xử lý không gian và thời gian. Sáu nhiệm vụ cụ thể:

\begin{description}[leftmargin=2.2em, style=nextline]
  \item[Nhiệm vụ 1 -- Khảo cứu lịch sử AI.] Hệ thống hóa các mốc phát triển từ thập niên 1940 đến nay, phân tích
    nguyên nhân hai “mùa đông AI”, sự chuyển dịch sang học máy thống kê, cuộc cách mạng học sâu và xu hướng hiện tại
    (mô hình nền tảng, RAG, hệ đa tác tử).
  \item[Nhiệm vụ 2 -- Cơ sở lý thuyết.] Trình bày hàm mất mát và cách tối ưu của từng họ mô hình: lan truyền ngược
    trong MLP, phép tích chập và pooling trong CNN, lan truyền ngược qua thời gian và cơ chế cổng trong RNN.
  \item[Nhiệm vụ 3 -- Khảo sát dữ liệu thật.] Phân tích khám phá (EDA) và tiền xử lý \textbf{bảy} tập dữ liệu thuộc
    năm lĩnh vực: y tế (\SL{eda/diabetes/n} hồ sơ tiểu đường), bất động sản Việt Nam (\SL{eda/house/n} tin đăng),
    thương mại điện tử (\SL{eda/reviews/n} nhận xét khách hàng), thị giác máy tính (SVHN
    \SL{eda/svhn/n_train}\,+\,\SL{eda/svhn/n_test} ảnh, EuroSAT \SL{eda/eurosat/n_train}\,+\,\SL{eda/eurosat/n_test} ảnh vệ tinh),
    tài chính (\SL{eda/stock/n_total} phiên giao dịch của bốn mã VN30) và hành vi người dùng
    (\SL{eda/churn/n_users} chuỗi nghe nhạc của KKBox).
  \item[Nhiệm vụ 4 -- Hiện thực trên ba nền tảng.] Cài đặt mô hình bằng \textbf{NumPy thuần} (tự viết lan truyền tiến,
    lan truyền ngược, tích chập, tế bào LSTM), bằng \textbf{Keras/TensorFlow} (API cấp cao) và bằng
    \textbf{PyTorch} (lớp \texttt{nn.Module}, vòng huấn luyện tự viết, chạy GPU).
  \item[Nhiệm vụ 5 -- Đánh giá và đối chuẩn.] So sánh các mô hình trên cùng một cách chia dữ liệu theo độ chính xác
    (Accuracy, Precision, Recall, F1, ROC-AUC, MAE, RMSE, $R^2$), số tham số, dung lượng tệp, thời gian huấn luyện
    và độ trễ suy luận đo trong cùng điều kiện phần cứng.
  \item[Nhiệm vụ 6 -- Triển khai.] Đóng gói các mô hình tốt nhất thành một ứng dụng web Flask gồm sáu phân hệ, nhận
    dữ liệu người dùng và trả kết quả dự báo theo thời gian thực.
\end{description}

\section*{3. Đối tượng và phương pháp nghiên cứu}

\paragraph{Đối tượng.} Các mô hình học máy cơ bản (Linear/Ridge Regression, Logistic Regression, Random Forest, SVM,
MLP); các kiến trúc CNN (Custom CNN, VGG-11, MobileNetV2, ResNet-18 và 1D CNN); các kiến trúc RNN (Simple RNN, LSTM,
GRU, BiLSTM); cùng bảy tập dữ liệu kể trên. Mọi mô hình học sâu được \emph{huấn luyện từ đầu} với trọng số ngẫu
nhiên, không dùng trọng số tiền huấn luyện, để kết quả phản ánh đúng năng lực của từng kiến trúc trên dữ liệu đích.

\paragraph{Phương pháp.} Kết hợp nghiên cứu lý thuyết với kiểm chứng thực nghiệm có kiểm soát:
\begin{itemize}
  \item cùng một hạt giống ngẫu nhiên (\texttt{seed = 42}) và cùng một cách chia huấn luyện/kiểm thử cho mọi mô hình
    trên một tập dữ liệu; bộ tiền xử lý chỉ được khớp trên phần huấn luyện;
  \item chia phân tầng theo nhãn cho bài toán phân loại mất cân bằng, chia theo trục thời gian (không xáo trộn) cho
    chuỗi giá chứng khoán để tránh rò rỉ thông tin tương lai;
  \item tập kiểm định tách riêng cho dừng sớm và chọn ngưỡng; tập kiểm thử chỉ dùng một lần để báo cáo;
  \item mỗi kết luận so sánh đi kèm một mốc tham chiếu (ví dụ dự báo “giữ nguyên giá hôm nay” cho chứng khoán);
  \item mọi con số trong báo cáo được sinh tự động từ tệp kết quả JSON của chính lần chạy, không chép tay.
\end{itemize}

\paragraph{Phân bổ phần cứng.} Máy thực nghiệm có CPU 32 luồng và GPU NVIDIA GeForce RTX 4060 Laptop (8 GB). PyTorch
dùng GPU; TensorFlow~2.21 trên Windows không hỗ trợ GPU nên các mô hình Keras chạy CPU. Các tiến trình độc lập chạy
song song: scikit-learn, mạng NumPy và Keras trên CPU (giới hạn số luồng mỗi tiến trình), PyTorch trên GPU.

\section*{4. Bố cục báo cáo}

\begin{itemize}
  \item \textbf{Chương 1 -- Lịch sử phát triển trí tuệ nhân tạo:} bức tranh hơn 80 năm từ nơ-ron hình thức đến hệ đa tác tử.
  \item \textbf{Chương 2 -- Các kỹ thuật học máy cơ bản:} lý thuyết bốn họ thuật toán, EDA ba tập dữ liệu bảng và văn bản,
    pipeline scikit-learn, MLP trên ba nền tảng, đối chuẩn và Web App.
  \item \textbf{Chương 3 -- Mạng nơ-ron tích chập:} phép tích chập, vùng cảm nhận, pooling, ba thế hệ kiến trúc;
    SVHN, EuroSAT và 1D CNN; hiện thực bằng NumPy, Keras và PyTorch; đối chuẩn độ chính xác -- chi phí; Web App thị giác.
  \item \textbf{Chương 4 -- Mạng nơ-ron hồi quy:} RNN, BPTT, triệt tiêu đạo hàm, cắt tỉa gradient, LSTM, GRU, BiLSTM;
    chứng khoán VN30 và KKBox; hiện thực ba nền tảng; đối chuẩn PyTorch -- Keras; Web App dự báo.
\end{itemize}
Phần Kết luận tổng hợp các phát hiện và hướng phát triển. Mã nguồn đầy đủ nằm trong thư mục
\texttt{src/TieuLuan 01} của kho lưu trữ; mỗi đoạn mã trong báo cáo được trích nguyên văn từ đó.
